% !TEX program = xelatex
\documentclass[12pt,a4paper]{ctexart}

% ========== 页面设置 ==========
\usepackage[top=2.5cm,bottom=2.5cm,left=2.5cm,right=2.5cm,headheight=15pt]{geometry}
\usepackage{amsmath,amssymb}
\usepackage{graphicx}
\usepackage{booktabs}
\usepackage{multirow}
\usepackage{array}
\usepackage{longtable}
\usepackage{float}
\usepackage{caption}
\usepackage{hyperref}
\usepackage{cite}
\usepackage{xcolor}
\usepackage{enumitem}
\usepackage{fancyhdr}
\usepackage{titlesec}
\setlength{\parskip}{6pt plus 2pt minus 1pt}
\setlength{\parindent}{0pt}

% ========== 页眉页脚 ==========
\pagestyle{fancy}
\fancyhf{}
\fancyhead[C]{中小微企业的信贷决策研究}
\fancyfoot[C]{\thepage}
\renewcommand{\headrulewidth}{0.4pt}

% ========== 标题格式 ==========
\titleformat{\section}{\Large\bfseries}{\thesection}{1em}{}
\titleformat{\subsection}{\large\bfseries}{\thesubsection}{1em}{}

% ========== 文档信息 ==========
\title{\vspace{-2cm}\heiti\zihao{2} 中小微企业的信贷决策研究\\
\vspace{0.3cm}\normalfont\zihao{4} ——基于AHP-熵权-TOPSIS的综合评价模型}
\author{\zihao{4} 冯浩然}
\date{}

\begin{document}
\maketitle

\begin{center}
\zihao{-5}
\vspace{0.5cm}
{\heiti 摘要：}
中小微企业信贷决策的核心难点在于如何在缺乏传统财务数据（资产负债表、现金流量表）的条件下，
仅凭交易票据信息量化企业信用风险。本文提出了一种\textbf{主客观融合}的信贷风险评估框架：
(1) 从进销项发票中提取15个企业级财务与行为指标，构建涵盖实力、稳定性、信誉、成长的四维评价体系；
(2) 采用AHP层次分析法确定主观权重（占60\%），熵权法确定客观权重（占40\%），通过组合赋权
兼顾专家经验与数据规律，并通过TOPSIS法获得企业信用评分；
(3) 利用附件1中的违约标签验证模型效果——风险等级从A到D的违约率呈严格单调递增（10.0\% $\to$
12.9\% $\to$ 16.2\% $\to$ 56.0\%），D级违约率是A级的5.6倍；
(4) 结合附件3数据建立利率-流失率二次关系（$R^2>0.99$），在总额约束下按风险得分差异化分配
信贷额度和利率；（5）引入行业冲击系数框架量化新冠疫情对不同行业企业的差异化影响。
灵敏度分析表明模型对权重扰动具有高度稳健性（Spearman $\rho=0.965$）。
通过三家企业（E38、E64、E115）的案例分析，验证了模型评分与实际经营状况的一致性。

\vspace{0.3cm}
{\heiti 关键词：}中小微企业；信贷风险；AHP层次分析法；熵权法；TOPSIS；利率定价；新冠疫情
\end{center}

\newpage
\tableofcontents
\newpage

% ============================================================
\section{问题重述}
% ============================================================

\subsection{原始问题}

某银行拟对中小微企业发放一年期贷款，额度为10--100万元，年利率范围为4\%--15\%。银行拥有以下数据：
\begin{itemize}
    \item \textbf{附件1}：123家有信贷记录企业的信息（含人工评定的A/B/C/D信誉评级、是否违约标签），
          以及2017--2020年的进项发票（210,947条）和销项发票（162,484条）；
    \item \textbf{附件2}：302家无信贷记录企业的信息，以及进项发票（395,175条）和销项发票（330,835条）；
    \item \textbf{附件3}：29组贷款利率与客户流失率的统计数据（按A/B/C三个信誉等级分别记录）。
\end{itemize}

需解决的三个问题：
\begin{enumerate}[label=(\arabic*)]
    \item 对附件1中123家企业的信贷风险进行量化分析，在\textbf{年度信贷总额固定}时给出信贷策略；
    \item 对附件2中302家企业的信贷风险进行量化分析，在\textbf{年度信贷总额为1亿元}时给出信贷策略；
    \item 考虑\textbf{突发因素}（如新冠疫情）对企业的影响，给出信贷调整策略。
\end{enumerate}

\subsection{问题的数学本质}

三个问题分别对应三种经典的决策分析范式：
\begin{itemize}
    \item \textbf{问题1}：\textbf{有监督的综合评价}。存在已知的违约标签可验证模型，核心是建立一个
          能有效区分违约与正常企业的评分函数 $f: \mathbb{R}^m \to [0,1]$；
    \item \textbf{问题2}：\textbf{无监督的外推评价}。将问题1中验证过的评价体系应用于无标签企业，
          核心是论证评价框架的\textbf{可迁移性}——即当企业特征空间的结构相似时，同一评分函数有效；
    \item \textbf{问题3}：\textbf{带外部冲击的鲁棒决策}。在问题2的基础上引入外生冲击向量
          $\mathbf{S} = (S_1, S_2, \ldots)$ 调整评分，核心是冲击系数的\textbf{合理量化}。
\end{itemize}

% ============================================================
\section{数据预处理与特征工程}
% ============================================================

\subsection{数据清洗}

将三个附件的Excel数据导入SQLite数据库。清洗工作包括：
\begin{enumerate}
    \item 删除附件2进项发票信息表中3列完全为空的数据列；
    \item 统一日期格式（附件1为日期，附件2为日期时间）；
    \item 规范化发票状态字段：附件1销项发票中有47条记录标记为``\ 作废发票''（含前导空格），
          统一修正为``作废发票''；
    \item 对\verb|企业代号|、\verb|开票日期|、\verb|发票状态|三列建立联合索引以加速后续聚合查询。
\end{enumerate}

\subsection{指标体系的构建逻辑}

在选择评价指标时，遵循\textbf{可获取性、可解释性、区分度}三个原则。信贷风险评估的核心问题是：
\textbf{这家企业未来一年能否按时还款？}发票数据虽然不直接包含资产负债表，但隐含着回答这个问题的关键信息：

\begin{itemize}
    \item \textbf{还款能力}：由经营规模和盈利能力决定——收入越大、利润越高，还款越有保障；
    \item \textbf{还款稳定性}：由经营波动决定——月度收入忽高忽低意味着现金流不稳定；
    \item \textbf{还款意愿}：由历史交易合规性决定——作废发票多、负数发票多可能意味着经营不规范；
    \item \textbf{发展前景}：由增长趋势决定——收入持续增长的企业更有可能维持还款能力。
\end{itemize}

基于上述逻辑，从发票数据中聚合计算了15个企业级指标，分为四个维度：

\begin{table}[H]
\centering
\caption{信贷风险评价指标体系}
\label{tab:indicators}
\begin{tabular}{c|l|l|c}
\toprule
\textbf{维度（准则）} & \textbf{指标} & \textbf{数据来源与计算方式} & \textbf{方向} \\
\midrule
\multirow{5}{*}{实力B1}
  & 销项总金额    & $\sum$(有效销项 $|$价税合计$|$)   & 极大型 \\
  & 进项总金额    & $\sum$(有效进项 $|$价税合计$|$)   & 极大型 \\
  & 利润          & 销项总金额 $-$ 进项总金额        & 极大型 \\
  & 利润率        & 利润/销项总金额 （截断至$[-0.5,0.5]$）& 极大型 \\
  & 交易网络广度  & 供应商数 $+$ 客户数            & 极大型 \\
\midrule
\multirow{2}{*}{稳定性B2}
  & 收入波动率    & $\sigma_{\text{月收入}}/\mu_{\text{月收入}}$ & 极小型 \\
  & 成本波动率    & $\sigma_{\text{月成本}}/\mu_{\text{月成本}}$ & 极小型 \\
\midrule
\multirow{4}{*}{信誉B3}
  & 销项有效率    & 有效发票数/总发票数            & 极大型 \\
  & 进项有效率    & 有效发票数/总发票数            & 极大型 \\
  & 销项作废率    & 作废发票数/总发票数            & 极小型 \\
  & 进项作废率    & 作废发票数/总发票数            & 极小型 \\
\midrule
\multirow{2}{*}{成长B4}
  & 收入环比增长率 & 最近两月收入变化率（截断至$[-1,2]$）& 极大型 \\
  & 综合活跃月数   & $\max$（进项月数, 销项月数）  & 极大型 \\
\bottomrule
\end{tabular}
\end{table}

\subsection{特征变换的必要性}

直接使用原始数据会导致模型被极端值主导。例如，某小型企业因上月收入接近零，
其``收入环比增长率''可达数十倍，远超正常范围。为此进行以下变换：

\begin{enumerate}
    \item \textbf{对数变换}：对金额类指标（销项总金额等）应用 $\ln(1+x)$。原因是企业规模
          呈幂律分布——少数大企业贡献大部分交易额，对数变换使模型不被少数超大企业主导；
    \item \textbf{非线性惩罚}：对作废率应用 $e^{3x}-1$。原因是作废率从5\%升到10\%的信用含义
          与从30\%升到35\%截然不同——后者意味着企业可能有系统性票据违规，需要指数级惩罚；
    \item \textbf{截断处理}：增长率限制在$[-1,2]$（即$-100\%$至$+200\%$），利润率限制在$[-0.5,0.5]$。
          超出范围的比率通常是``小分母问题''（如新开企业月收入极小导致增长率异常），
          而非真实的经营信息。
\end{enumerate}

% ============================================================
\section{模型选择：为什么是AHP+熵权+TOPSIS？}
% ============================================================

在构建评价模型前，需要回答三个设计选择问题。

\subsection{赋权方法的选择：为什么用组合赋权？}

常用的赋权方法可分为三类：

\begin{table}[H]
\centering
\caption{赋权方法对比}
\label{tab:weight_comparison}
\begin{tabular}{c|l|l|l}
\toprule
\textbf{方法} & \textbf{类型} & \textbf{优点} & \textbf{缺点} \\
\midrule
AHP & 主观 & 融入专家经验，可解释 & 依赖专家判断质量 \\
熵权法 & 客观 & 基于数据，无主观偏差 & 可能被噪声指标误导 \\
CRITIC法 & 客观 & 考虑指标冲突性 & 计算复杂，不适合大规模 \\
\bottomrule
\end{tabular}
\end{table}

\textbf{选择组合赋权（AHP + 熵权法）而非单一方法的理由}：

纯熵权法存在一个致命问题：如果某个指标本身波动大但与信用风险无关（例如``收入环比增长率''受季节性
和开业时间影响极大），它会被赋予过高权重。在附件1的实际计算中，收入环比增长率的原始熵权高达
0.388——但这并不意味着它比其他指标重要4倍。AHP通过专家经验对此进行了纠正（AHP权重仅0.048），
组合后降至0.184。

\textbf{选择AHP而非纯熵权法的理由}：信贷风险评价是领域知识密集型问题——银行信贷员知道
``作废发票多''比``规模小''更危险。这种先验知识必须通过AHP等主观赋权方法引入模型。

\textbf{选择熵权法而非CRITIC法的理由}：CRITIC法需要计算指标间相关系数矩阵，在15个指标且数据
分布偏态的情况下，Pearson相关系数不够稳健；而熵权法基于信息熵，对分布形态不敏感。

\subsection{组合比例的选择：为什么$\alpha=0.6$？}

组合权重公式为 $w = \alpha w_{\text{AHP}} + (1-\alpha) w_{\text{Entropy}}$。$\alpha$的取值
需要在``相信专家''和``相信数据''之间权衡。

本文选择$\alpha=0.6$的依据：
\begin{enumerate}
    \item 对$\alpha \in \{0.3, 0.4, 0.5, 0.6, 0.7\}$进行了违约分离度测试（违约企业与正常企业
          的TOPSIS均值差），$\alpha=0.6$时分离度最大（差值0.052）；
    \item $\alpha=0.6$时，收入环比增长率组合权重从0.388降至0.184，但仍保留了一定影响力——
          成长性确实对信用有预测作用，但不应主导模型；
    \item 在信贷领域，专家的判断通常比纯数据更可靠，因为数据可能反映相关性而非因果性。
\end{enumerate}

\subsection{评价方法的选择：为什么用TOPSIS？}

常用的多指标综合评价方法包括TOPSIS、灰色关联分析（GRA）和模糊综合评价（FCE）。本文选择TOPSIS
的理由如下：

\begin{table}[H]
\centering
\caption{评价方法适用性对比}
\label{tab:method_comparison}
\begin{tabular}{c|l|l}
\toprule
\textbf{方法} & \textbf{核心思想} & \textbf{不适用于本题的原因} \\
\midrule
GRA & 与最优参考序列的关联度 & 需要设定参考序列；样本数大时区分度弱 \\
FCE & 模糊集合的隶属度 & 依赖评语集的主观划分；不适合连续数值指标 \\
\textbf{TOPSIS} & \textbf{距理想解的欧氏距离} & \textbf{直观、非线性、适合大规模评价} \\
\bottomrule
\end{tabular}
\end{table}

TOPSIS的优势在于：（1）利用欧氏距离的非线性性质，对异常值有自然的``拉近''效应——即一个指标
极差但其他指标极好的企业不会被简单地平均；（2）结果在$[0,1]$上有自然的概率解释；
（3）不需要预先设定评语集或参考序列，减少了主观干扰环节。

% ============================================================
\section{信贷风险评估模型}
% ============================================================

\subsection{AHP层次分析法}

AHP的核心是将复杂的多指标评价问题分解为层次结构，通过两两比较确定各指标的相对重要性。
本文建立三层结构：\textbf{目标层（信贷风险） $\to$ 准则层（B1实力/B2稳定性/B3信誉/B4成长）
$\to$ 指标层（15个指标）}。

\subsubsection{准则层判断矩阵的构造依据}

准则层两两比较的标度基于以下信贷风控常识：
\begin{itemize}
    \item \textbf{实力 vs 稳定性（3:1）}：一家收入高但波动大的公司，银行通常仍愿放贷（只是利率更高）；
          但一家收入稳定却极低的公司，银行很难放贷。实力的重要性显著高于稳定性；
    \item \textbf{实力 vs 信誉（2:1）}：银行更看重``还不起''（能力）而非``不想还''（意愿），
          因为后者可通过法律手段约束。但差距不宜过大——严重违规（高作废率）也是真实风险；
    \item \textbf{信誉 vs 稳定性（2:1）}：交易合规性是比经营波动更强的风险信号；
    \item \textbf{成长 vs 实力（1:4）}：成长性是``锦上添花''，实力才是``雪中送炭''。
\end{itemize}

据此构造准则层判断矩阵：

\begin{equation}
\mathbf{B} = \begin{bmatrix}
1   & 3   & 2   & 4 \\
1/3 & 1   & 1/2 & 2 \\
1/2 & 2   & 1   & 3 \\
1/4 & 1/2 & 1/3 & 1
\end{bmatrix}
\end{equation}

采用几何平均法计算权重：
\begin{equation}
w_k = \frac{\left(\prod_{j=1}^{n} b_{kj}\right)^{1/n}}{\sum_{i=1}^{n}\left(\prod_{j=1}^{n} b_{ij}\right)^{1/n}}
\end{equation}
得到 $\mathbf{w}_B = (0.467, 0.160, 0.278, 0.095)$。

一致性检验：$\lambda_{\max} = 4.031$，$CI = 0.011$，$CR = CI/RI(4) = 0.011/0.89 = 0.0093 < 0.10$，
通过检验。这意味着判断矩阵的逻辑自洽性良好——不存在``A比B重要，B比C重要，但C又比A重要''的逻辑矛盾。

\subsection{熵权法}

熵权法基于信息论原理：如果所有企业在某个指标上的取值几乎相同，那么该指标对区分企业信用好坏
几乎没有贡献，应赋予低权重；反之，差异大的指标信息量大，应赋予高权重。

计算步骤如下：
\begin{enumerate}
    \item \textbf{比重矩阵}：$p_{ij} = x'_{ij} / \sum_{i=1}^{n} x'_{ij}$，其中$x'_{ij}$为正向化后的值；
    \item \textbf{信息熵}：$e_j = -\frac{1}{\ln n}\sum_{i=1}^{n} p_{ij} \ln(p_{ij} + \varepsilon)$，
          $\varepsilon = 10^{-12}$防止$\ln(0)$；
    \item \textbf{差异系数}：$d_j = 1 - e_j$；
    \item \textbf{权重}：$w_j^{\text{Ent}} = d_j / \sum d_j$。
\end{enumerate}

附件1的熵权法结果显示，收入环比增长率的熵权（0.388）远高于其他指标，这是因为部分新开业企业
的增长率极高，造成数据离散度大。但如前文模型选择部分所述，这并不意味着增长率的重要性是交易
有效率的15倍。这正是组合赋权$\alpha=0.6$需要发挥作用的地方。

\subsection{组合赋权与TOPSIS评分}

组合权重为 $w_j = 0.6 \cdot w_j^{\text{AHP}} + 0.4 \cdot w_j^{\text{Ent}}$。
TOPSIS计算步骤为：正向化 $\to$ 向量归一化 $\to$ 加权矩阵 $\to$ 正负理想解 $\to$ 欧氏距离
$\to$ 相对贴近度 $C_i = D_i^-/(D_i^++D_i^-)$。$C_i$越大表示企业越优，按降序排列得风险排名。

\subsection{模型验证}

验证模型的标准是：它能否将已知的27家违约企业和96家正常企业有效区分？

\subsubsection{均值差异检验}

\begin{itemize}
    \item 27家违约企业的平均TOPSIS得分为0.181；
    \item 96家正常企业的平均TOPSIS得分为0.233；
    \item 差值0.052。这在实际意义上意味着\textbf{模型将违约企业平均排在正常企业之后}。
\end{itemize}

\subsubsection{风险等级与违约率的单调性}

将123家企业按TOPSIS得分从高到低分为A（前25\%）、B（25\%--50\%）、C（50\%--80\%）、
D（后20\%）四档。这是模型最重要的验证——好的模型应该让风险等级与违约率严格单调：

\begin{table}[H]
\centering
\caption{风险等级违约率分布（附件1验证）}
\label{tab:default_validation}
\begin{tabular}{c|c|c|c|c}
\toprule
\textbf{等级} & \textbf{企业数} & \textbf{违约数} & \textbf{违约率} & \textbf{含义} \\
\midrule
A & 30 & 3 & 10.0\% & 模型认为最安全的群体，违约率最低 \\
B & 31 & 4 & 12.9\% & 较安全，违约率略高于A \\
C & 37 & 6 & 16.2\% & 中等风险，违约率进一步上升 \\
D & 25 & 14 & \textbf{56.0\%} & 模型认为最危险的群体，违约率剧烈攀升 \\
\bottomrule
\end{tabular}
\end{table}

违约率从A级的10.0\%到D级的56.0\%呈严格单调递增，D级违约率是A级的5.6倍。
这一结果验证了模型的\textbf{序数有效性}——虽然具体概率值可能因样本有限而波动，
但模型给出的\textbf{排名顺序}与真实风险高度一致。

\subsubsection{模型评级与银行人工评级的差异分析}

\begin{table}[H]
\centering
\caption{模型评级 vs 银行评级交叉表}
\label{tab:crosstab_detail}
\begin{tabular}{c|cccc|c}
\toprule
& \textbf{模型A} & \textbf{模型B} & \textbf{模型C} & \textbf{模型D} & \textbf{银行合计} \\
\midrule
银行A & 10 & 11 & 6  & 0  & 27 \\
银行B & 9  & 9  & 17 & 3  & 38 \\
银行C & 8  & 8  & 9  & 9  & 34 \\
银行D & 1  & 4  & 6  & 13 & 24 \\
\midrule
模型合计 & 28 & 32 & 38 & 25 & 123 \\
\bottomrule
\end{tabular}
\end{table}

两种评级的完全一致率仅为33.3\%，但\textbf{相邻一致率（相差不超过1级）达到81.3\%}。
这不一致是正常的，且恰恰说明了模型的独特价值：

\begin{itemize}
    \item 银行给E103评为D级的理由可能不在发票数据中（如企业主个人信用差、涉及诉讼等），
          但该企业发票数据表现优秀（利润率60\%，有效发票率超84\%），模型基于客观数据给出了较高评分。
          \textbf{这不是模型错误，而是模型补充了银行可能忽略的正面信息}。
    \item 银行给E64评为A级，但模型给出了C级。检查发现该企业收入波动极大（月度变异系数>2）、
          利润率微薄（$<$5\%）。模型捕捉到了发票数据中隐含的风险信号。
\end{itemize}

% ============================================================
\section{典型企业案例分析}
% ============================================================

为了验证模型评分与企业实际经营面貌的一致性，选取三家代表性企业进行深入分析。

\subsection{案例1：E38 — 模型评分最高的企业（0.848，A级）}

\begin{table}[H]
\centering
\caption{E38 关键指标与数据集均值对比}
\label{tab:case1}
\begin{tabular}{l|c|c|c}
\toprule
\textbf{指标} & \textbf{E38的值} & \textbf{附件1均值} & \textbf{解读} \\
\midrule
销项总金额 & 数亿元级 & 约1.33亿 & 远超平均，属于头部企业 \\
利润率 & $>$20\% & 约48\% & 盈利健康 \\
销项客户数（对数） & 高 & 230 & 客户网络广泛 \\
收入波动率 & $<$0.5 & 0.91 & 远低于平均，经营稳定 \\
销项有效率 & $>$99\% & 88\% & 几乎无作废发票，交易规范性极高 \\
销项作废率 & $<$1\% & 12\% & 极少发票异常 \\
收入环比增长率 & 温和正增长 & 0.86 & 稳定增长，无异常波动 \\
\bottomrule
\end{tabular}
\end{table}

E38在所有维度上均表现出众：规模大、利润高、经营稳定、发票规范、适度增长。银行原始评级为B级，
模型给出的A级评分结合了客观数据提供的更全面信息。该企业应获得最高额度（100万）和最低利率（4\%）。

\subsection{案例2：E64 — 银行A级但模型C级（0.188，C级）}

\begin{table}[H]
\centering
\caption{E64 关键指标分析}
\label{tab:case2}
\begin{tabular}{l|c|c|l}
\toprule
\textbf{指标} & \textbf{E64的值} & \textbf{均值} & \textbf{风险解读} \\
\midrule
销项总金额 & 中等偏上 & 1.33亿 & 实力尚可 \\
利润率 & 极低 & 48\% & 盈利能力弱——赚得多但留得少 \\
收入波动率 & $>$2.0 & 0.91 & 月度收入大起大落，现金流不可预测 \\
成本波动率 & $>$3.0 & 1.22 & 成本控制不稳定 \\
销项有效率 & 约75\% & 88\% & 四分之一发票存在问题 \\
收入环比增长率 & $-0.95$ & 0.86 & 最近两月收入断崖式下跌 \\
\bottomrule
\end{tabular}
\end{table}

E64虽然在绝对规模上不差，但呈现出\textbf{多重预警信号}：收入剧烈波动 + 成本失控 + 发票有效率低 +
近期收入暴跌。银行可能基于历史关系或行业地位给了A级，但当前发票数据揭示的风险值得银行重新审视。
建议给该企业C级对待：中等额度、偏高利率，同时密切监控其经营恢复情况。

\subsection{案例3：E115 — 模型D级且确实违约（0.139，D级，已违约）}

\begin{table}[H]
\centering
\caption{E115 关键指标分析}
\label{tab:case3}
\begin{tabular}{l|c|c|l}
\toprule
\textbf{指标} & \textbf{E115的值} & \textbf{均值} & \textbf{风险解读} \\
\midrule
销项总金额 & 极低（$<$百万级） & 1.33亿 & 规模极小 \\
利润 & 负数 & 约3789万 & 入不敷出 \\
利润率 & $-100\%$ & 48\% & 持续亏损 \\
销项作废率 & $>$15\% & 12\% & 高作废率，交易规范性差 \\
收入波动率 & $>$1.5 & 0.91 & 不稳定 \\
收入环比增长率 & 负 & 0.86 & 持续下滑 \\
\bottomrule
\end{tabular}
\end{table}

E115几乎在所有维度上都处于最差水平：规模极小、持续亏损、发票不规范、经营下滑。
该企业被模型评为D级且确实为违约企业，验证了模型对\textbf{全面恶化型企业}的有效识别。

\textbf{三个案例的启示}：模型评分并非简单的规模排序，而是对多个风险维度的综合权衡。
E64（规模较大但多重预警）得分低于一些规模更小但经营稳健的企业，说明模型能够捕捉超越规模的
深层风险信号。

% ============================================================
\section{信贷策略模型}
% ============================================================

\subsection{利率-流失率关系}

附件3给出了29组利率-流失率对应数据。对其分别拟合二次多项式：

\begin{equation}
L_k(r) = a_k r^2 + b_k r + c_k, \quad k \in \{A,B,C\}
\end{equation}

拟合结果：A级 $R^2=0.993$，B级 $R^2=0.995$，C级 $R^2=0.995$，均接近完美拟合。三条曲线具有
相同的经济含义：\textbf{流失率随利率加速上升}（二次系数为负值表明曲线凹向原点，低利率区间流失率
增长缓慢，高利率区间流失率急剧攀升）。这为差异化定价提供了定量依据——对A级企业低利率可以
``锁定''优质客户（流失率低），对C级企业高利率虽然会增加流失率但利率溢价带来的收入增量可能
超过流失损失。

\subsection{信贷额度与利率分配}

\textbf{额度分配}遵循三步：
\begin{enumerate}
    \item \textbf{筛除}：D级企业不放贷（题目规定）；
    \item \textbf{比例法}：$L_i = \frac{C_i}{\sum_{j\in\mathcal{E}} C_j} \times \text{Budget}$；
    \item \textbf{剪裁+再分配}：限制在$[10,100]$万，通过迭代再分配使总额收敛到预算（误差$<0.05\%$）。
\end{enumerate}

\textbf{利率定价}采用``等级基准+得分微调''双层结构：
\begin{enumerate}
    \item \textbf{第一层：按风险等级定基准}。基准利率的设定综合考虑了利率下限（4\%）、
          流失率曲线（低利率锁定优质客户）、以及不同等级企业的风险溢价；
    \item \textbf{第二层：同等级内按得分微调}。$r_i = r_{\text{base}} + (0.5 - \text{score\_norm}_i) \times 0.02$，
          最高得分企业获得最低利率（利率下限4\%），最低得分企业利率上浮1\%。
\end{enumerate}

双层结构的意义在于：\textbf{等级间的利率跳跃反映质的差异（A vs C本质不同），等级内的微小调整
反映量的差异（同为A级，E38比E89更优）}。

\subsection{问题1与问题2的信贷策略}

问题1中``年度信贷总额固定''但未给定具体数字。本文基于以下推理确定总额为\textbf{8,000万元}：
\begin{itemize}
    \item 123家企业，排除D级25家后剩余98家；
    \item 按10--100万约束，中位数约80万时恰好充分利用预算；
    \item $98\times 80 = 7,840 \approx 8,000$ 万元。这是一个使约束紧凑（高分企业接近100万、
          低分企业接近10万）的合理数字。
\end{itemize}

\begin{table}[H]
\centering
\caption{两个问题的信贷策略对比}
\label{tab:strategy_comparison}
\begin{tabular}{l|c|c}
\toprule
\textbf{指标} & \textbf{问题1（123家，8000万）} & \textbf{问题2（302家，1亿）} \\
\midrule
放贷企业数（占比） & 98家（79.7\%）  & 241家（79.8\%） \\
不放贷企业数（D级）& 25家            & 61家 \\
额度范围 & 64--100万      & 31--100万 \\
额度中位数 & 81万           & 37万 \\
满额100万企业 & 15家          & 9家 \\
A级平均利率 & 4.60\%         & 4.68\% \\
B级平均利率 & 7.38\%         & 7.43\% \\
C级平均利率 & 10.96\%        & 10.97\% \\
预期利息收入 & $\approx$570万/年 & $\approx$790万/年 \\
预算利用率 & 100.0\%         & 100.0\% \\
\bottomrule
\end{tabular}
\end{table}

问题2的额度中位数（37万）显著低于问题1（81万），原因是302家企业的得分分布更加分散——
大量小微企业得分较低导致分配更均匀，而非集中于少数几家头部企业。

% ============================================================
\section{突发因素下的信贷调整}
% ============================================================

\subsection{行业识别}

从企业名称中提取行业关键词进行分类。由于企业名称经过脱敏处理（如``***科技发展有限公司''），
仅约38\%的企业可精确识别行业。分类结果如下：
\begin{itemize}
    \item 严重负面（旅游/酒店/餐饮/影院）：3家；
    \item 中度负面（建筑/制造/批发零售/物流）：69家；
    \item 正面/轻影响（医疗/科技/电商/IT）：43家；
    \item 无法分类（其他）：187家。
\end{itemize}

\subsection{冲击系数量化}

冲击系数的设定需要回答一个问题：\textbf{疫情对不同行业企业的一年期违约概率影响有多大？}
基于清华大学与北京大学对995家中小企业的联合调研以及《清华金融评论》的数据\cite{ref_covid_liu}：

\begin{itemize}
    \item 旅游/餐饮/酒店业：2020年春节黄金周收入损失超过\textbf{5,000亿元}，行业平均营收预计下降
          \textbf{50\%以上}，约\textbf{85\%}的企业现金储备不足以支撑3个月。据此估计违约概率上升
          约30--50\%，取保守值30\%，对应冲击系数 $S=1.30$；
    \item 制造/建筑/批发零售业：受供应链中断和消费萎缩双重影响，约\textbf{58\%}的中小企业预计
          2020年营收降幅超过50\%。考虑到这些行业有一定资产可变现，估计违约概率上升约15\%，
          对应 $S=1.15$；
    \item 医疗/科技/电商行业：在疫情期间需求反而增长。医疗行业订单量上升、在线服务需求激增。
          估计违约概率下降约10\%，对应 $S=0.90$。
\end{itemize}

冲击系数的经济含义是：$S_k = \frac{\text{调整后违约概率}}{\text{调整前违约概率}}$。
TOPSIS得分除以冲击系数即可获得调整后的风险排序：$C'_i = C_i / S_k$。

\subsection{调整结果}

共83家企业（27.5\%）的风险等级发生了变化。严重负面行业的3家企业平均贷款额度从40.1万骤降至
22.4万（下降44\%），利率从5.72\%升至约6.8\%；正面行业的43家企业平均额度从30.7万升至37.5万
（上升22\%），利率获得约1\%的优惠。调整方向符合政策直觉：\textbf{保护受冲击企业的最有效方式
不是增加贷款（可能变成坏账），而是将有限的信贷资源向抗风险能力强的行业倾斜}。

% ============================================================
\section{灵敏度分析与鲁棒性}
% ============================================================

\subsection{权重扰动}

对每个指标权重进行100次随机$\pm 10\%$独立扰动，比较扰动后企业排名与原始排名的秩相关性：
Spearman相关系数均值为0.965，Kendall相关系数均值为0.848。Spearman > 0.85说明模型排名
对权重扰动具有高度稳健性——即使某个指标的权重被错误设定10\%，最终排名也不会发生剧烈变化。

\subsection{指标的独立贡献}

逐一移除每个指标后重新计算排名。移除收入环比增长率后Spearman相关系数最低（0.026），
这反映了该指标对排名的影响最大——但它同时也说明组合赋权（$\alpha=0.6$）起到了预期的
``牵制''作用：如果完全采用熵权法（$\alpha=0$），该指标将完全主导排名。

% ============================================================
\section{模型的优点与局限}
% ============================================================

\subsection{优点}

\begin{enumerate}
    \item \textbf{主客观融合的赋权策略}：AHP（60\%）+ 熵权（40\%）的组合避免了纯主观的随意性
          和纯客观被噪声误导的双重风险，$\alpha=0.6$的选择经过了违约分离度检验；
    \item \textbf{可解释的指标体系}：每个指标都对应明确的信贷风控含义（能力/稳定性/意愿/前景），
          且经过合理的数学变换（对数压缩、指数惩罚、截断处理）；
    \item \textbf{严格单调的违约率验证}：模型给出的风险排序与实际违约率完全单调，D级违约率56.0\%
          是A级的5.6倍，证明模型具有真实的违约识别能力；
    \item \textbf{有依据的冲击量化}：疫情冲击系数基于公开的学术调研数据而非凭空设定；
    \item \textbf{已验证的鲁棒性}：Spearman $\rho=0.965$ 证明模型对权重设定不敏感。
\end{enumerate}

\subsection{局限与改进方向}

\begin{enumerate}
    \item \textbf{时序信息的利用不充分}：当前模型仅使用了聚合统计量和环比增长率，未利用发票的
          完整时间序列。例如，收入``先升后降''与``先降后升''在聚合指标中可能看起来相同，
          但信用含义截然不同。引入时间序列特征（如趋势斜率、拐点检测）可能提升预测精度；
    \item \textbf{行业分类精度不足}：62\%的企业名称无法识别行业。可考虑利用发票中的交易对手信息
          进行行业推断——例如供应商和客户的行业分布可以侧面反映企业所处产业链位置；
    \item \textbf{信贷策略未使用优化模型}：当前采用规则化的比例分配+分级定价。更精确的方法
          是将利率和额度作为连续决策变量，构建以银行期望收益最大化为目标的非线性规划模型，
          用启发式算法（粒子群、遗传算法）求解；
    \item \textbf{未利用上下游网络信息}：发票数据本质上是企业间的交易图谱。图神经网络（GNN）
          可捕捉``交易对手的对手违约了，我的风险也会上升''的传染效应。
\end{enumerate}

% ============================================================
\section{结论}
% ============================================================

本文从中小微企业信贷决策的实际需求出发，基于发票数据构建了完整的风险量化与信贷策略框架。
核心贡献在于：（1）建立了\textbf{主客观融合}（AHP+熵权）的赋权方法和TOPSIS评价模型，
经验证能有效区分违约与正常企业；（2）通过三家典型企业的案例分析，展示了模型评分与
实际经营面貌的一致性；（3）为三个问题分别给出了\textbf{数据驱动、逻辑自洽、定量可验}
的信贷策略方案。

在总额约束下，D级企业不放贷是基于风险控制的审慎原则；A/B/C级企业的差异化额度与利率
是基于``高风险高利率、低风险低利率''的市场化定价原则；疫情调整策略则是基于
``将有限资源向抗风险能力强的企业倾斜''的危机管理原则。

模型的工程实现使用了SQLite处理百万级发票数据，Python科学计算生态（NumPy/Pandas/Matplotlib）
完成建模与可视化，完整代码开源托管于：\\
\url{https://github.com/wanbro333/evalution_model_test}

\begin{thebibliography}{99}

\bibitem{saaty} Saaty T L. \textit{The Analytic Hierarchy Process}. McGraw-Hill, 1980.

\bibitem{hwang} Hwang C L, Yoon K. \textit{Multiple Attribute Decision Making}. Springer, 1981.

\bibitem{shannon} Shannon C E. A mathematical theory of communication. \textit{Bell System Technical Journal}, 1948, 27(3): 379--423.

\bibitem{cheng} 程扬, 周大勇, 程帆等. 基于组合赋权法的中小微企业信贷风险量化及预测. \textit{系统工程}, 2023, 41(1): 140--151.

\bibitem{luo} 罗胜尹, 彭睿杰, 胡静. 基于中小微企业信贷风险和银行收益的信贷决策研究. \textit{中国集体经济}, 2021(27): 79--80.

\bibitem{ref_covid_liu} 刘超, 王海军. 新冠肺炎疫情对中小企业的冲击与金融纾困. \textit{清华金融评论}, 2020(6).

\bibitem{ref_covid_cba} 中国银行业协会. 新冠肺炎疫情下的中小微企业经营风险分析和应对. 2020年4月.

\bibitem{ref_tsinghua} 清华大学商业模式创新研究中心, 北京大学汇丰商学院. 995家中小企业受疫情影响调研报告. 2020年2月.

\end{thebibliography}

\end{document}
